\section{Topic names}\label{sec:session_names}

A \emph{topic name} is a string of printable ASCII characters used by the application
to identify a publish-subscribe channel.
The session layer carries names only in gossip and scout messages; every other
message identifies topics by their 64-bit hash (section~\ref{sec:session_identity}).

\subsection{Character set and normalization}

A valid name consists of characters in the ASCII range $[33, 126]$; the space
character (ASCII $32$) and all control characters are excluded.
Names are case-sensitive.
The forward-slash character \texttt{/} (ASCII $47$) acts as the \emph{separator}
between name segments.

A \emph{normalized} name contains no leading separator, no trailing separator, and
no consecutive separators. Names are normalized before hashing and before being sent
on the wire. The maximum length of a normalized name is 200 bytes.

\subsection{Pattern tokens}

Subscribers may register under pattern names that contain the following special
tokens; publishers may not. A publisher whose resolved name contains any pattern
token shall be rejected by the session layer.

\begin{description}
    \item[\texttt{*} (ASCII $42$)] Matches exactly one name segment.
    For example, the pattern \texttt{sensors/*/data} matches \texttt{sensors/imu/data}
    and \texttt{sensors/gps/data} but not \texttt{sensors/imu/raw/data}.

    \item[\texttt{>} (ASCII $62$)] Matches zero or more trailing segments.
    Must appear as the last token in the pattern.
    For example, the pattern \texttt{sensors/>} matches \texttt{sensors},
    \texttt{sensors/imu}, and \texttt{sensors/imu/raw/data}.
\end{description}

Pattern tokens are not permitted in intermediate positions outside of these forms;
for example, \texttt{sen*ors} is a literal name, not a pattern.

\subsection{Pin suffix}

A name may be pinned to a specific subject-ID using the suffix
\texttt{\#\textit{NNN}}, where \texttt{\#} is ASCII $35$ and \textit{NNN} is a
decimal number in the range $[0, 8191]$ without leading zeros.
Pinning is permitted only on verbatim names; a pattern name shall not carry a
pin suffix.
The pin suffix is stripped from the name before hashing and normalization, so a
pinned topic and its verbatim counterpart share the same topic hash but request a
different subject-ID (section~\ref{sec:session_identity}).

\subsection{Home and namespace expansion}

A name may begin with the \emph{home} token \texttt{\textasciitilde} (tilde,
ASCII $126$), which is replaced at resolution time with the node's home name.
A name beginning with \texttt{/} is \emph{absolute} and is used as-is, ignoring
the namespace prefix. All other names are relative and are prefixed with the
node's namespace at resolution time.

\subsection{Remapping}

A node may carry a \emph{remap table} of \texttt{\textit{from}}$\to$\texttt{\textit{to}}
pairs that rewrite names before namespace and home expansion.
A matching rule replaces the input name with its \texttt{\textit{to}} string and then
feeds the result into the rest of the resolution pipeline.
If the matched rule's \texttt{\textit{to}} string carries a pin suffix, that pin is
honoured; any pin supplied in the user's original name is discarded.
Software implementations should accept the remap table from the environment variables
\texttt{CYPHAL\_NAMESPACE} and \texttt{CYPHAL\_REMAP} to allow integrators to adjust
names without modifying the application.

\subsection{Resolution pipeline}

The following pipeline is normative. Given a user-supplied name, a node resolves it
as follows, in order:

\begin{enumerate}
    \item \textbf{Remap} the input name through the remap table, if any.
    \item \textbf{Expand} the home token \texttt{\textasciitilde} if present, or
    prefix the namespace if the name is relative and the absolute form is not used.
    \item \textbf{Extract} any pin suffix, recording the pinned subject-ID separately.
    \item \textbf{Normalize} the remaining name by collapsing consecutive separators
    and trimming leading and trailing separators.
\end{enumerate}

The result is the normalized resolved name (used for hashing and for inclusion in
gossip and scout messages) together with the optional pinned subject-ID.
Resolution fails if the normalized name is empty, exceeds 200 bytes, or contains a
character outside the permitted set.

\begin{remark}
    The intended usage pattern is that applications hardcode verbatim names and let
    integrators introduce patterns and pinning through the remap table at deployment
    time. This keeps the application code free of site-specific routing decisions.
\end{remark}

\subsection{Resolution examples}

The following examples illustrate the resolution pipeline when no remap table
is installed. \texttt{NAMESPACE} and \texttt{HOME} are per-node configuration;
\texttt{RESOLVED} is the normalized name that appears on the wire;
\texttt{PIN} is the pinned subject-ID, if any;
\texttt{VERBATIM?} is \emph{yes} if the resolved name has no pattern tokens.

\begin{CyphalSimpleTable}[wide]{Name resolution without remap}{|l l l l r l|}
    Name                   & Namespace   & Home    & Resolved           & Pin    & Verbatim? \\
    \texttt{foo/bar}       & \texttt{ns1}   & \texttt{me} & \texttt{ns1/foo/bar}       & --       & yes \\
    \texttt{\textasciitilde/foo/bar} & \texttt{ns1} & \texttt{me} & \texttt{me/foo/bar} & -- & yes \\
    \texttt{/foo//bar/}    & \texttt{ns1}   & \texttt{me} & \texttt{foo/bar}           & --       & yes \\
    \texttt{foo/bar/}      & \texttt{\textasciitilde/ns1} & \texttt{me} & \texttt{me/ns1/foo/bar} & -- & yes \\
    \texttt{foo\#123}      & \texttt{ns1\#456} & \texttt{me} & \texttt{ns1/foo}         & $123$    & yes \\
    \texttt{foo/\#123}     & \texttt{ns1\#456} & \texttt{me} & \texttt{ns1/foo}         & $123$    & yes \\
    \texttt{*/foo}         & \texttt{ns1}   & \texttt{me} & \texttt{ns1/*/foo}         & --       & no  \\
    \texttt{\textasciitilde/*/foo/} & \texttt{ns1} & \texttt{me} & \texttt{me/*/foo}    & --       & no  \\
    \texttt{/\textasciitilde/*/foo/} & \texttt{ns1} & \texttt{me} & \texttt{\textasciitilde/*/foo} & -- & no \\
\end{CyphalSimpleTable}

Note in the last row that an absolute name beginning with \texttt{/}
short-circuits home expansion, so the tilde is preserved as a literal
character in the resolved name.

Next, a remap table is installed. A matched remap rule replaces the user-supplied
name with the rule's \texttt{to} string, which is then fed through the same
pipeline; a matched rule also discards any pin suffix the user may have
supplied. The pin on a rule's \texttt{to} string, if present, is honoured
instead.

\begin{CyphalSimpleTable}[wide]{Name resolution with remap}{|l l l l l r X|}
    Name                 & From           & To                & Namespace & Home    & Pin    & Remark \\
    \texttt{foo/bar}     & \texttt{foo/bar} & \texttt{zoo}     & \texttt{ns} & \texttt{me} & --   & Relative remap; resolved name is \texttt{ns/zoo}. \\
    \texttt{foo/bar}     & \texttt{foo/bar} & \texttt{zoo\#123} & \texttt{ns} & \texttt{me} & $123$ & Pinned relative remap; resolved name is \texttt{ns/zoo}. \\
    \texttt{foo/bar\#456} & \texttt{foo/bar} & \texttt{zoo}     & \texttt{ns} & \texttt{me} & --   & User pin discarded by the matched rule; resolved name is \texttt{ns/zoo}. \\
    \texttt{foo/bar}     & \texttt{foo/bar} & \texttt{/zoo}    & \texttt{ns} & \texttt{me} & --   & Absolute remap; namespace is ignored; resolved name is \texttt{zoo}. \\
    \texttt{foo/bar}     & \texttt{foo/bar} & \texttt{\textasciitilde/zoo} & \texttt{ns} & \texttt{me} & -- & Homeful remap; the tilde expands to \texttt{me}; resolved name is \texttt{me/zoo}. \\
\end{CyphalSimpleTable}

The following inputs cause resolution to fail:

\begin{description}
    \item[\texttt{"foo bar\textbackslash{}nbaz"}] Contains a space (ASCII $32$)
    and a control character (ASCII $10$); non-printable and space characters
    are not allowed.
    \item[\texttt{"foo/*/bar\#123"}] A pattern name cannot be pinned; the
    \texttt{\#} suffix is not permitted on names containing pattern tokens.
    \item[\texttt{""} (empty)] An empty name is not allowed.
    \item[\texttt{"\#1234"}] After the pin expression is stripped, the
    remaining name is empty, which is not allowed.
    \item[A name longer than 200 bytes after normalization] Rejected regardless
    of how much buffer space is available.
\end{description}

\section{Topic identity and subject-ID mapping}\label{sec:session_identity}

Every topic has a deterministic 64-bit \emph{topic hash} derived from its normalized
resolved name and a transport-visible \emph{subject-ID} derived from the hash and
its eviction counter.
The mapping is purely arithmetic and uses no shared state other than the parameters
normalized in this section; every node computes the same value from the same inputs.

\subsection{The topic hash function}\label{sec:session_rapidhash}

The topic hash is computed as $\mathrm{rapidhash}(\mathit{name},\ \mathit{seed}=0)$,
where \textit{name} is the UTF-8 byte sequence of the normalized resolved name
(no NUL terminator) and $\mathrm{rapidhash}$ is the 64-bit
\textbf{rapidhash V3} function\footnote{%
    rapidhash is a fast, high-quality, platform-independent hash function
    developed by Nicolas De Carli and distributed under the MIT license at
    \href{https://github.com/Nicoshev/rapidhash}{github.com/Nicoshev/rapidhash}.
    It is derived from the wyhash algorithm by Wang Yi.
    Cyphal v1.1 pins the V3 variant of rapidhash as its normative hash function;
    the Python transcription below has been verified bit-identical to the
    reference C implementation for every test vector in this section.
} with the default secret specified below.

The default secret consists of eight 64-bit constants:

\begin{CyphalSimpleTable}{Default secret for the rapidhash V3 topic hash function}{|l l|}
    Index   & Value (hexadecimal)      \\
    0       & \texttt{0x2d358dccaa6c78a5} \\
    1       & \texttt{0x8bb84b93962eacc9} \\
    2       & \texttt{0x4b33a62ed433d4a3} \\
    3       & \texttt{0x4d5a2da51de1aa47} \\
    4       & \texttt{0xa0761d6478bd642f} \\
    5       & \texttt{0xe7037ed1a0b428db} \\
    6       & \texttt{0x90ed1765281c388c} \\
    7       & \texttt{0xaaaaaaaaaaaaaaaa} \\
\end{CyphalSimpleTable}

Multi-byte integers within the hash function are read from the input in
\emph{little-endian} byte order on every platform; big-endian implementations shall
emulate this ordering. Wherever the resulting 64-bit hash integer appears in a
session-layer header, it is serialized as little-endian.

\begin{remark}
    Multiple rapidhash variants exist in the wild
    (V1, V2, V3, V3-micro, V3-nano, wyhash, \ldots) and they are not wire-compatible
    with one another. Cyphal v1.1 pins down the full V3 variant only, with the
    default seed and the default secret listed above.
\end{remark}

A normative Python transcription of the reference C implementation is given in
appendix~\ref{sec:appendix_rapidhash}. It is bit-identical to the reference
implementation on every platform and may be used directly to reproduce the
test vectors below.

The following test vectors are authoritative. An implementation that fails to
reproduce any of them is not interoperable with the rest of the Cyphal network and
cannot participate in the topic allocation protocol.

\begin{CyphalSimpleTable}{Rapidhash V3 test vectors}{|l r l|}
    Input (ASCII)                            & Length & \texttt{rapidhash(input)}  \\
    \texttt{""} (empty)                      & 0      & \texttt{0x0338dc4be2cecdae} \\
    \texttt{"a"}                             & 1      & \texttt{0x599f47df33a2e1eb} \\
    \texttt{"abc"}                           & 3      & \texttt{0xcb475beafa9c0da2} \\
    \texttt{"topic"}                         & 5      & \texttt{0x0164a16804e84989} \\
    \texttt{"sensors/imu"}                   & 11     & \texttt{0xbb402fd6c693c60b} \\
    \texttt{"sensors/engine/temperature"}    & 26     & \texttt{0x32d2c4bd5c5485e7} \\
    \texttt{"head/camera/upper\_left"}       & 22     & \texttt{0x28f52efde90dec0a} \\
    \texttt{"system/primary/sensor/0/data"}  & 28     & \texttt{0x57627b4b782b47af} \\
    \texttt{"ins/0/data"}                    & 10     & \texttt{0xdf3339ab99373b98} \\
    \texttt{"1234"}                          & 4      & \texttt{0xa5ae810dbac9be9b} \\
    \texttt{"65535"}                         & 5      & \texttt{0x2c5a542e44fa7c20} \\
    200 $\times$ \texttt{'x'}                & 200    & \texttt{0x25adb114ca9ea76a} \\
\end{CyphalSimpleTable}

\subsection{Subject-ID mapping}\label{sec:session_subject_id_mapping}

The subject-ID assigned to a non-pinned topic is a function of its hash $h$, its
eviction counter $e$, and the transport's subject-ID modulus $M$:
\begin{equation*}
    \mathit{subject\_id}(h, e) = 8192 + ((h + e \cdot e) \bmod M)
\end{equation*}
All arithmetic is performed on 64-bit unsigned integers with wraparound;
implementations using signed types shall take care not to sign-extend the hash.
The modulus $M$ is a prime satisfying $M \bmod 4 = 3$, chosen per
section~\ref{sec:session_modulus}; its selection ensures that the
$\{e^2 \bmod M : e \in \mathbb{N}\}$ sequence is a valid quadratic-probing
sequence.

A \emph{pinned} topic is one whose eviction counter lies in the reserved range
$[\text{UINT32\_MAX}-8191,\ \text{UINT32\_MAX}]$.
Its subject-ID is obtained by $\mathit{subject\_id} = \text{UINT32\_MAX} - e$,
which lies in the pinned range $[0, 8191]$.
Pinned topics otherwise behave as ordinary CRDT records and participate in
arbitration on equal terms with non-pinned topics.

\begin{remark}
    The high-eviction encoding of pinned topics serves a second purpose: it acts as a
    wrap-around guard for the eviction counter itself.
    Under any pathological churn that drove a non-pinned topic's eviction counter
    past $2^{32} - 8192$, the topic would enter the pinned range and be arrested at
    subject-ID 8191 rather than overflowing the counter and disrupting consensus.
    The 32-bit counter is wide enough that such growth is not reachable in practice:
    at a sustained one eviction per second it takes over 136 years to overflow.
\end{remark}

\subsection{Subject-ID ranges}\label{sec:session_subject_id_ranges}

The session layer partitions the entire subject-ID space exposed by the
transport into four disjoint ranges, described in the table below.
All values refer to the numerical subject-ID carried in the transport header;
the topic $\to$ subject-ID mapping is defined in
section~\ref{sec:session_subject_id_mapping}.

\begin{CyphalSimpleTable}{Subject-ID space partition}{|l X|}
    Range                          & Usage \\
    $[0,\ 8191]$                   & Pinned topics. Retained as the primary path for interoperation with Cyphal/CAN v1.0 nodes, which only understand this range. \\
    $[8192,\ 8191{+}M]$            & Auto-allocated topics. $M$ is a transport-specific prime modulus satisfying $M \bmod 4 = 3$, defined in section~\ref{sec:session_modulus}. \\
    $(8191{+}M,\ B)$               & Gossip shard subjects used for background CRDT dissemination (section~\ref{sec:session_gossip}). \\
    $B$                            & The \emph{broadcast subject}, used for broadcast gossip and for scout service discovery (section~\ref{sec:session_scout}). \\
\end{CyphalSimpleTable}

where $B = 2^{\lceil\log_2(8191 + M + 1)\rceil} - 1$ is the highest subject-ID
representable in the transport's subject-ID width.
Every node is expected to subscribe to the broadcast subject; the transport's
implicit participant discovery mechanism (section~\ref{sec:transport}) is a
side effect of that universal subscription.

On the broadcast subject and on any gossip shard subject, the transport layer
is permitted to relax two of the requirements stated in
section~\ref{sec:transport}: deduplication of incoming transfers is not
mandatory, and support for transfers exceeding a single transport frame is not
required. The session layer tolerates occasional duplication on these
subjects, and constrains its own protocol messages on them to fit within a
single transport frame.

\begin{remark}
    The relaxation is important for scalability on resource-constrained
    transports. Deduplication state on the broadcast subject would otherwise
    grow linearly with the number of peers, which is hostile to small nodes;
    multi-frame reassembly buffers on the gossip shards would impose a
    per-shard memory cost that a large network cannot afford.
\end{remark}

\subsection{Subject-ID modulus}\label{sec:session_modulus}

The modulus $M$ depends on the width of the subject-ID field supported by the
transport. It shall be prime, satisfy $M \bmod 4 = 3$, and be the largest such
prime not exceeding $(\text{transport maximum subject-ID}) - 8191$.
The following values are normative:

\begin{CyphalSimpleTable}{Subject-ID moduli for supported transport widths}{|l l r X|}
    Name    & Width   & Modulus $M$     & Notes                                                             \\
    16-bit  & 16 bits & $57203$         & Usable on all Cyphal transports, including Cyphal/CAN.             \\
    23-bit  & 23 bits & $8378431$       & Usable on Cyphal/UDP; not compatible with Cyphal/CAN.              \\
    32-bit  & 32 bits & $4294954663$    & Reserved for future 32-bit transports; not compatible with the above. \\
\end{CyphalSimpleTable}

All nodes on the same logical Cyphal network shall use the same modulus.
In a heterogeneously redundant deployment where different transports are used in
parallel, the smallest applicable modulus shall be chosen so that every subject-ID
is representable on every member transport.

For the 16-bit modulus the derived subject-ID ranges are $[0, 8191]$ for pinned
topics, $[8192, 65394]$ for auto-allocated topics, $[65395, 65534]$ for the
140 gossip shards, and $65535$ for the broadcast subject.
For the 23-bit modulus the ranges are $[0, 8191]$, $[8192, 8386622]$,
$[8386623, 8388606]$ ($1984$ gossip shards), and $8388607$ respectively.

The following test vectors lock down the mapping formula for typical inputs.
The evictions column traces the first three entries of the quadratic-probing
sequence $(e = 0, 1, 2)$; subsequent entries step by $5, 7, 9, \ldots$.

\begin{CyphalSimpleTable}{Subject-ID test vectors}{|l l r r r|}
    Name               & Hash                 & $e$ & $M=57203$ & $M=8378431$ \\
    \texttt{""}        & \texttt{0x0338dc4be2cecdae} & 0 & $21600$  & $6520572$ \\
    \texttt{""}        & \texttt{0x0338dc4be2cecdae} & 1 & $21601$  & $6520573$ \\
    \texttt{""}        & \texttt{0x0338dc4be2cecdae} & 2 & $21604$  & $6520576$ \\
    \texttt{"topic"}   & \texttt{0x0164a16804e84989} & 0 & $38082$  & $1049799$ \\
    \texttt{"topic"}   & \texttt{0x0164a16804e84989} & 1 & $38083$  & $1049800$ \\
    \texttt{"topic"}   & \texttt{0x0164a16804e84989} & 2 & $38086$  & $1049803$ \\
    \texttt{"sensors/imu"} & \texttt{0xbb402fd6c693c60b} & 0 & $54214$ & $8263034$ \\
    \texttt{"sensors/imu"} & \texttt{0xbb402fd6c693c60b} & 1 & $54215$ & $8263035$ \\
    \texttt{"sensors/imu"} & \texttt{0xbb402fd6c693c60b} & 2 & $54218$ & $8263038$ \\
    \texttt{"ins/0/data"} & \texttt{0xdf3339ab99373b98} & 0 & $28748$ & $8266023$ \\
    \texttt{"ins/0/data"} & \texttt{0xdf3339ab99373b98} & 1 & $28749$ & $8266024$ \\
    \texttt{"ins/0/data"} & \texttt{0xdf3339ab99373b98} & 2 & $28752$ & $8266027$ \\
\end{CyphalSimpleTable}
